\subsection{Derivatives and Integrals}\label{sec:distribution-derivatives-integrals}

\subsubsection{Basic Definitions}\label{subsec:distribution-derivative-basic-definitions}

In this section, we discuss differentiation of distributions and various applications. We shall confine our discussion to distributions in $\mathcal{D}^{\prime}(\Omega)$; completely analogous considerations apply in $\mathcal{S}^{\prime}(\mathbb{R}^m)$. We define the derivative of a distribution as follows.

\begin{Definition}
\label{def:distribution-partial-derivative}
Let $f\in\mathcal{D}^{\prime}(\Omega)$. Then the \textbf{derivative} of $f$ with respect to $x_j$ is defined as
\begin{equation}\label{eq:distribution-partial-derivative-definition}
\left(\dfrac{\partial f}{\partial x_j},\phi\right)=-\left(f,\dfrac{\partial\phi}{\partial x_j}\right).
\end{equation}

\end{Definition}

\begin{Remark}
\label{rem:distribution-classical-derivative-agreement}
If $f$ is in $C^1(\Omega)$, this definition agrees with the classical derivative, as can be seen by an integration by parts. It is easy to see that $\partial f/\partial x_j$ is again in $\mathcal{D}^{\prime}(\Omega)$.
\end{Remark}

Note that differentiation is a continuous operation in $\mathcal{D}^{\prime}(\Omega)$, i.e., the reader can show the following.

\begin{Theorem}
\label{thm:distribution-derivative-continuity}
If $f_n\to f$ in $\mathcal{D}^{\prime}(\Omega)$ then $\partial f_n/\partial x_j\to\partial f/\partial x_j$.
\end{Theorem}

Thus, for distributions exchanging derivatives and limits is no problem, quite a contrast to the situation in classical calculus.

\begin{Remark}
\label{rem:distribution-higher-derivatives}
Higher derivatives are defined recursively; the equality of mixed partial derivatives is obvious from the definition and the equality of the mixed partial derivatives of test functions. In general, we have
\begin{equation}\label{eq:distribution-higher-derivative-pairing}
(D^\alpha f,\phi)=(-1)^{|\alpha|}(f,D^\alpha\phi).
\end{equation}

\end{Remark}

\begin{Remark}
\label{rem:distribution-difference-quotients}
The classical derivative is defined as a limit of difference quotients. In a sense, distributional derivatives are still limits of difference quotients. In the previous section, we defined the translation of a distribution by
\begin{equation}\label{eq:distribution-coordinate-translation-pairing}
(f(\bm{x}+h\bm{e}_j),\phi(\bm{x}))=(f(\bm{x}),\phi(\bm{x}-h\bm{e}_j)).
\end{equation}
This does not necessarily make sense, because $\bm{x}-h\bm{e}_j$ need not lie in $\Omega$. For fixed $\phi\in\mathcal{D}(\Omega)$, however, $\phi(\bm{x}-h\bm{e}_j)$ is in $\mathcal{D}(\Omega)$ provided $h$ is sufficiently small. Hence \eqref{eq:distribution-coordinate-translation-pairing} is meaningful for small $h$, although how small $h$ has to be depends on $\phi$. We now find\[
\begin{aligned}
&\lim_{h\to0}\dfrac1h\bigl[(f(\bm{x}+h\bm{e}_j),\phi(\bm{x}))-(f(\bm{x}),\phi(\bm{x}))\bigr]\\
&\quad=\lim_{h\to0}\left(f(\bm{x}),\dfrac1h(\phi(\bm{x}-h\bm{e}_j)-\phi(\bm{x}))\right)\\
&\quad=\left(f(\bm{x}),\lim_{h\to0}\dfrac1h(\phi(\bm{x}-h\bm{e}_j)-\phi(\bm{x}))\right)\\
&\quad=-\left(f,\dfrac{\partial\phi}{\partial x_j}\right)\\
&\quad=\left(\dfrac{\partial f}{\partial x_j},\phi\right).
\end{aligned}
\]
\end{Remark}
